\section{模型的建立与求解}\label{sec:model-establishment-solution}

\subsection{总体分析}\label{sec:model-overview}

四个问题构成由定位数据处理到任务决策的递进过程。问题一针对无噪声异频数据完成时间对齐与轨迹重建；问题二在随机噪声和固定空间差条件下进行时空标定与异步融合；问题三根据实际数据的偏差检验结果选择融合状态；问题四利用所得固定轨迹安排射击和拍照任务。各问题的输入与输出见表\ref{tab:model-interfaces}。

统一采用方式1时钟作为参考时钟，位置、时间分别以米和秒为单位，轨迹输出频率为10 Hz。时间对齐结果用于构造两设备的同步残差，融合后的位移、速度和加速度用于计算任务距离与完整准备区间内的约束裕度。

\begin{table}[htbp]
\centering
\caption{四个问题的输入输出关系}
\label{tab:model-interfaces}
\footnotesize
\setlength{\tabcolsep}{3pt}
\begin{tabular}{c c p{2.5cm} p{3.2cm} p{3.3cm}}
\toprule
问题 & 主要输入 & 处理内容 & 核心输出 & 后续用途 \\
\midrule
问题一 & 附件1无噪声定位 & 时间对齐与插值重建 & 时间校正量、10 Hz位置 & 确定时间映射形式 \\
问题二 & 附件2含噪定位 & 稳健标定与异步融合 & 时空偏差、10 Hz融合状态 & 验证含噪融合方法 \\
问题三 & 附件3实际定位 & 偏差检验与状态选择 & 偏差判定、状态与标准差 & 提供任务计算轨迹 \\
问题四 & 问题三轨迹、附件4目标 & 可行窗口与任务调度 & 九项日程与约束裕度 & 填写结果工作簿副本 \\
\bottomrule
\end{tabular}
\end{table}

总体技术路线见图\ref{fig:technical-route}。原始4 Hz和5 Hz观测先完成时间配准，再根据数据条件进入无噪声重建或含噪异步融合；偏差检验决定问题三是否需要增设偏差状态，所得轨迹用于固定路径上的任务优化。

\begin{figure}[htbp]
  \centering
  \includegraphics[width=.86\textwidth]{paper/technical_route.pdf}
  \caption{多源定位、状态融合与任务调度的总体技术路线}
  \label{fig:technical-route}
\end{figure}

各问题之间的计算关系见图\ref{fig:subproblem-flow}。问题三输出位置、速度、加速度及模型内位置标准差；问题四据此判断准备区间的连续可行性和任务冲突，不重新规划机器人路径。

\begin{figure}[htbp]
  \centering
  \includegraphics[width=.88\textwidth]{paper/subproblem_flow.pdf}
  \caption{四个问题的计算关系与共享状态量}
  \label{fig:subproblem-flow}
\end{figure}

\subsection{模型一的建立与求解：无噪声时间配准与10 Hz轨迹重建}\label{sec:model-q1}

\subsubsection{模型一的建立}\label{sec:q1-establishment}

\noindent\textbf{（1）模型目标、输入输出与证据边界。}

附件 1 给出同一机器人轨迹的两组无噪声采样：方式 1 为 4 Hz，方式 2 为 5
Hz。两台设备开机先后不同，故设备时间轴之间存在固定偏差。本问需要估计该
偏差，并在不外推的前提下生成 10 Hz 二维位置轨迹。

数据审计表明，方式 1、方式 2 分别有 3000、4001 个有效观测，采样间隔严格
为 0.25 s、0.20
s，无缺失和重复时间戳。由于题面明确无噪声，本问不进行滤波、
平滑或异常值删除。

\noindent\textbf{（2）方法选择与无噪声边界。}

方案 A 直接在粗网格上最小化位置误差，简单但最终精度受网格步长限制；方案
B
先用异频统一网格上的速度互相关产生候选，再以全局位置误差扫描和有界连续优化
精估时间偏差，可同时防止周期轨迹的伪相关峰；方案 C
采用动态时间规整允许非线性
时间伸缩，但题面只存在固定开机偏差，额外自由度会改变问题含义。因此选择方案
B。

连续轨迹比较 Linear、Cubic Spline、PCHIP 和 Akima
四种插值。最终以对齐后的 跨传感器位置 RMSE
为主判据，并兼顾无噪声轨迹的连续可导性，选择 Cubic Spline；其中 PCHIP 的保形性质及其使用边界参照分段三次插值理论\cite{fritsch1980}。

\noindent\textbf{（3）时间映射与位置误差目标。}

记两种定位数据为

\begin{equation}
D_m=\{(t_i^{(m)},x_i^{(m)},y_i^{(m)})\},\qquad m=1,2,
\label{eq:q1-data}
\end{equation}

并分别建立连续二维轨迹 \(\mathbf r_m(t)\)。定义方式 2 的校正时间为

\begin{equation}
t_{2,\mathrm{aligned}}=t_2+\Delta t.
\label{eq:q1-time-map}
\end{equation}

因此 \(\Delta t<0\) 表示方式 2 的时间戳需要整体向前平移。对候选偏差
\(\Delta t\)，公共有效区间为

\begin{equation}
\Omega(\Delta t)=
\left[
\max(t_{1,\min},t_{2,\min}+\Delta t),
\min(t_{1,\max},t_{2,\max}+\Delta t)
\right].
\label{eq:q1-overlap}
\end{equation}

在公共评价时刻 \(\{t_k\}\) 上定义位置均方误差

\begin{equation}
J(\Delta t)=\frac{1}{N}
\sum_{k=1}^{N}
\left\|\mathbf r_1(t_k)-\mathbf r_2(t_k-\Delta t)\right\|_2^2,
\label{eq:q1-objective}
\end{equation}

从而得到

\begin{equation}
\Delta t^*=\arg\min_{\Delta t}J(\Delta t).
\label{eq:q1-optimum}
\end{equation}

\subsubsection{模型一的求解}\label{sec:q1-solution}

求解时先在 0.1 s 公共网格上计算二维速度特征互相关，保留多个相关峰；随后在
全部可行偏差范围内以 0.5 s 步长扫描位置
MSE，最后在最优候选附近用有界标量
优化精细求解。精细优化使用固定评价网格，避免候选偏差改变公共区间起点所造成的
数值非光滑。

\noindent\textbf{（1）同时刻观测合并与10 Hz重建。}

对方式 2 时间戳施加 \(\Delta t^*\)
后，将两组观测合并、排序。若两个观测落在
同一真实时刻，则先检查坐标差，再取均值保留一个节点。对合并后的 6301
个唯一 时刻建立 Cubic Spline，并在

\begin{equation}
t_n=221.0+0.1n,\qquad n=0,1,\ldots,8494
\label{eq:q1-grid}
\end{equation}

上输出 10 Hz 轨迹。输出区间为 \([221.0,1070.4]\)
s，仅覆盖至少一种原始观测 支持的时间区间，不进行外推。

图\ref{fig:q1-objective}表明周期轨迹会形成多个局部低谷，因此互相关只用于候选生成，最终偏差必须由式\eqref{eq:q1-objective}在全可行域内确定。

\begin{figure}[htbp]
  \centering
  \includegraphics[width=.88\textwidth]{q1/objective_scan.pdf}
  \caption{模型一全可行域时间偏差目标函数扫描}
  \label{fig:q1-objective}
\end{figure}

\subsubsection{算法实现}\label{sec:q1-algorithm}

\begin{pseudocodeblock}{算法1\quad 无噪声时间配准与10 Hz轨迹重建}
{\heiti 输入：}两路无噪声定位序列 \(D_1,D_2\)，目标输出频率10 Hz。\par
{\heiti 输出：}时间校正量 \(\Delta t^*\) 与10 Hz二维位置轨迹。\par
\begin{enumerate}[label={\heiti 步骤~\arabic*:},leftmargin=5.5em,itemsep=1pt,topsep=3pt]
  \item 分别构造两路连续轨迹，由速度特征互相关生成时间偏差候选集；
  \item 按式\eqref{eq:q1-overlap}确定可行域，在固定公共网格上计算式\eqref{eq:q1-objective}；
  \item 取全域扫描的最低谷，在其邻域内连续优化，得到式\eqref{eq:q1-optimum}的数值解；
  \item 按式\eqref{eq:q1-time-map}校正方式2时间戳；若两路观测对应同一时刻，则核验坐标后合并；
  \item 对全部唯一时刻建立三次样条，按式\eqref{eq:q1-grid}生成10 Hz轨迹并返回结果。
\end{enumerate}
\end{pseudocodeblock}

\subsubsection{结果分析}\label{sec:q1-results}

采用符号约定
\(t_{2,\mathrm{aligned}}=t_2+\Delta t\)，连续优化得到原始估计

\begin{equation}
\Delta t^*=-198.4316999986\ \mathrm{s}.
\label{eq:q1-raw-offset}
\end{equation}

考虑输入时间戳的显示精度，正式报告为

\begin{equation}
\boxed{\Delta t^*=-198.4317\ \mathrm{s}}.
\label{eq:q1-reported-offset}
\end{equation}

即方式 2 的时间轴需整体向前平移 198.4317 s。对齐公共区间长 700.35 s，
位置 RMSE 为 \(8.51\times10^{-11}\) m。最终 10 Hz 结果共 8495 行。

图\ref{fig:q1-aligned}显示校正后两路观测在公共区间内重合；图\ref{fig:q1-trajectory}给出在原始观测支持域内重建的完整10 Hz轨迹。两图共同说明式\eqref{eq:q1-time-map}只调整时间坐标，没有改变任何原始位置观测。

\begin{figure}[htbp]
  \centering
  \includegraphics[width=.88\textwidth]{q1/aligned_trajectory.pdf}
  \caption{模型一时间校正后的两路定位轨迹（公共区间700.35 s）}
  \label{fig:q1-aligned}
\end{figure}
\begin{figure}[htbp]
  \centering
  \includegraphics[width=.84\textwidth]{q1/trajectory_10hz.pdf}
  \caption{模型一在221.0--1070.4 s内重建的10 Hz完整轨迹}
  \label{fig:q1-trajectory}
\end{figure}

\subsubsection{模型检验与稳健性}\label{sec:q1-validation}

四种插值模型均给出相同的报告偏差。其对齐 RMSE 分别为：Linear
\(4.76\times10^{-6}\) m、Cubic Spline \(8.51\times10^{-11}\) m、PCHIP
\(2.53\times10^{-7}\) m、Akima \(9.91\times10^{-9}\) m，Cubic Spline
最优。

正确偏差下共有 700 组两设备共同采样时刻，最大坐标不一致量为
0，构成不依赖 插值评价网格的独立核验。合成数据测试中，预设偏差
\(-3.7317\) s 的恢复误差为 \(2.20\times10^{-10}\) s。

进一步改变最小重叠时长、粗搜索步长和精细评价步长，共运行 8
个预设情景；原始 时间偏差估计的极差仅为 \(3.56\times10^{-9}\)
s，所有情景按输入精度均报告 \(-198.4317\) s。最优点距搜索边界至少 100.35
s，排除了边界伪最优。

四种插值的误差数量级比较见图\ref{fig:q1-model-comparison}。图中四种方法给出的时间偏差一致，而三次样条的跨设备位置RMSE最低，说明方法选择由无噪声轨迹的局部曲率表达能力决定，并未改变时间偏差结论。

\noindent\textbf{适用边界。}

本模型依赖``固定时间偏差、无空间系统偏差、无随机噪声''三个题面条件。若设备时钟
存在漂移或测量含噪，直接插值会低估不确定性，因此问题二、三需改用偏差辨识和
统计融合。问题一输出的符号约定、公共区间处理和 10 Hz
数据契约将作为后续各问的统一基准。

\begin{figure}[htbp]
  \centering
  \includegraphics[width=.78\textwidth]{paper/model_comparison.pdf}
  \caption{模型一四种插值方法的跨设备位置RMSE比较}
  \label{fig:q1-model-comparison}
\end{figure}

\subsection{模型二的建立与求解：稳健时空标定与异步融合}\label{sec:model-q2}

\subsubsection{模型二的建立}\label{sec:q2-establishment}

\noindent\textbf{（1）含噪观测模型与时间映射。}

设机器人在参考时间轴上的真实连续轨迹为
\(\mathbf p^\star(t)=(x^\star(t),y^\star(t))^\mathsf T\)。以定位方式 1
的时钟为参考，两种定位方式的离散观测满足

\begin{equation}
\begin{aligned}
\mathbf z_i^{(1)}
  &=\mathbf p^\star\!\left(t_i^{(1)}\right)
    +\boldsymbol\varepsilon_i^{(1)},\\
\mathbf z_j^{(2)}
  &=\mathbf p^\star\!\left(t_j^{(2)}-\Delta t\right)
    +\mathbf b+\boldsymbol\varepsilon_j^{(2)},
\end{aligned}
\label{eq:q2-observation}
\end{equation}

其中，\(\Delta t\) 为方式 2 相对于方式 1 的固定时间偏差，
\(\mathbf b=(b_x,b_y)^\mathsf T\) 为方式 2 相对于方式 1
的固定空间偏差，随机误差满足

\begin{equation}
\mathbb E\!\left(\boldsymbol\varepsilon^{(m)}\right)=\mathbf 0,\qquad
\operatorname{Cov}\!\left(\boldsymbol\varepsilon^{(m)}\right)=\mathbf R_m,
\qquad m=1,2.
\label{eq:q2-noise-assumption}
\end{equation}

由式\eqref{eq:q2-observation}，方式 2 的设备时刻映射到参考时间轴为

\begin{equation}
t_{2,\mathrm{ref}}=t_2-\Delta t.
\label{eq:q2-time-map}
\end{equation}

分别由两组离散观测构造连续位置函数
\(\widetilde{\mathbf z}_1(t)\) 与
\(\widetilde{\mathbf z}_2(t)\)。对候选时间偏差
\(\delta\in\mathcal D\)，其公共支持区间为

\begin{equation}
\mathcal I(\delta)=
[t_{1,\min},t_{1,\max}]
\cap[t_{2,\min}-\delta,t_{2,\max}-\delta].
\label{eq:q2-overlap}
\end{equation}

为使不同候选值在相同样本集合上评价，取固定公共区间

\begin{equation}
\mathcal I_{\mathrm{fix}}
=\bigcap_{\delta\in\mathcal D}\mathcal I(\delta),
\qquad
\mathcal T_{\mathrm{fix}}=\{t_k\}_{k=1}^{N}
\subset\mathcal I_{\mathrm{fix}}.
\label{eq:q2-fixed-grid}
\end{equation}

\noindent\textbf{（2）时间偏差的粗估计。}

对连续位置函数求导，得到两种定位方式的速度序列
\(\mathbf v_m(t)=\mathrm d\widetilde{\mathbf z}_m(t)/\mathrm dt\)。
在粗网格上定义二维速度特征的归一化互相关

\begin{equation}
C(h)=
\frac{\displaystyle
\sum_{t\in\mathcal T_h}
\left(\mathbf v_1(t)-\overline{\mathbf v}_1\right)^\mathsf T
\left(\mathbf v_2(t+h)-\overline{\mathbf v}_2\right)}
{\displaystyle
\sqrt{\sum_{t\in\mathcal T_h}
\left\|\mathbf v_1(t)-\overline{\mathbf v}_1\right\|_2^2}
\sqrt{\sum_{t\in\mathcal T_h}
\left\|\mathbf v_2(t+h)-\overline{\mathbf v}_2\right\|_2^2}},
\label{eq:q2-ncc}
\end{equation}

其中，\(\mathcal T_h\) 为平移后仍位于两组观测支持域内的粗评价网格。
粗时间偏差取为

\begin{equation}
h_c=\arg\max_{h\in\mathcal H} C(h),
\label{eq:q2-coarse-offset}
\end{equation}

并据此在 \(h_c\) 邻域内构造精搜索域 \(\mathcal D\)。

\noindent\textbf{（3）Huber 剖面标定。}

对任意候选 \(\delta\)，在固定公共网格上构造两设备的位置差

\begin{equation}
\mathbf d_k(\delta)=
\widetilde{\mathbf z}_2(t_k+\delta)
-\widetilde{\mathbf z}_1(t_k).
\label{eq:q2-residual}
\end{equation}

对坐标分量 \(r\in\{x,y\}\)，采用 MAD 给出固定稳健尺度

\begin{equation}
s_r=1.4826\,
\operatorname{median}_k
\left|d_{kr}(\delta_0)
-\operatorname{median}_{\ell}d_{\ell r}(\delta_0)\right|,
\label{eq:q2-robust-scale}
\end{equation}

其中 \(\delta_0\) 为精搜索域中点。尺度在整个精搜索过程中保持不变。
令

\begin{equation}
u_{kr}=\frac{d_{kr}(\delta)-b_r}{s_r},\qquad
\rho_c(u)=
\begin{cases}
\dfrac{u^2}{2},& |u|\le c,\\[2mm]
c|u|-\dfrac{c^2}{2},& |u|>c,
\end{cases}
\label{eq:q2-huber-loss}
\end{equation}

则给定 \(\delta\) 时，空间偏差由逐轴 Huber 估计确定：

\begin{equation}
\widehat b_r(\delta)
=\arg\min_{b_r}
\sum_{k=1}^{N}s_r^2
\rho_c\!\left(
\frac{d_{kr}(\delta)-b_r}{s_r}
\right).
\label{eq:q2-profile-bias}
\end{equation}

采用 IRLS 求解式\eqref{eq:q2-profile-bias}。第 \(\ell\) 次迭代的权重与更新量为

\begin{equation}
w_{kr}^{(\ell)}=
\min\!\left(1,\frac{c}{|u_{kr}^{(\ell)}|}\right),
\qquad
b_r^{(\ell+1)}
=\frac{\sum_k w_{kr}^{(\ell)}d_{kr}(\delta)}
{\sum_k w_{kr}^{(\ell)}}.
\label{eq:q2-irls-bias}
\end{equation}

将 \(\widehat{\mathbf b}(\delta)\) 回代，得到一维剖面目标

\begin{equation}
J_H(\delta)=
\frac{2}{N}\sum_{k=1}^{N}\sum_{r\in\{x,y\}}
s_r^2\rho_c\!\left(
\frac{d_{kr}(\delta)-\widehat b_r(\delta)}{s_r}
\right),
\label{eq:q2-huber-objective}
\end{equation}

从而

\begin{equation}
\widehat{\Delta t}
=\arg\min_{\delta\in\mathcal D}J_H(\delta),
\qquad
\widehat{\mathbf b}
=\widehat{\mathbf b}(\widehat{\Delta t}).
\label{eq:q2-calibration-estimator}
\end{equation}

\noindent\textbf{（4）测量噪声协方差估计。}

局部常加速度运动的三阶差分为零。对等间隔位置观测定义

\begin{equation}
\Delta^3\mathbf z_i^{(m)}
=\mathbf z_{i+3}^{(m)}-3\mathbf z_{i+2}^{(m)}
+3\mathbf z_{i+1}^{(m)}-\mathbf z_i^{(m)}.
\label{eq:q2-third-difference}
\end{equation}

若相邻测量误差独立同分布，则每个坐标分量满足
\(\operatorname{Var}(\Delta^3 z_r)=20\sigma_r^2\)。因此

\begin{equation}
\widehat\sigma_{mr}
=\frac{1.4826}{\sqrt{20}}\,
\operatorname{median}_i
\left|
\Delta^3z_{ir}^{(m)}
-\operatorname{median}_{\ell}\Delta^3z_{\ell r}^{(m)}
\right|,
\qquad
\widehat{\mathbf R}_m
=\operatorname{diag}
(\widehat\sigma_{mx}^2,\widehat\sigma_{my}^2).
\label{eq:q2-measurement-covariance}
\end{equation}

\noindent\textbf{（5）异步常加速度状态空间模型。}

将状态写为

\begin{equation}
\mathbf x(t)=
\bigl(\mathbf p(t)^\mathsf T,
\mathbf v(t)^\mathsf T,
\mathbf a(t)^\mathsf T\bigr)^\mathsf T
=(x,y,v_x,v_y,a_x,a_y)^\mathsf T.
\label{eq:q2-state}
\end{equation}

相邻两个观测事件的时间间隔为 \(\tau_k\) 时，状态转移模型为

\begin{equation}
\mathbf x_k=\mathbf F(\tau_k)\mathbf x_{k-1}+\mathbf w_k,
\qquad
\mathbf F(\tau)=
\begin{bmatrix}
\mathbf I_2&\tau\mathbf I_2&\dfrac{\tau^2}{2}\mathbf I_2\\
\mathbf 0&\mathbf I_2&\tau\mathbf I_2\\
\mathbf 0&\mathbf 0&\mathbf I_2
\end{bmatrix}.
\label{eq:q2-transition}
\end{equation}

采用白噪声 jerk 模型，其离散过程噪声协方差为

\begin{equation}
\mathbf Q(\tau;q)=q
\begin{bmatrix}
\dfrac{\tau^5}{20}\mathbf I_2&
\dfrac{\tau^4}{8}\mathbf I_2&
\dfrac{\tau^3}{6}\mathbf I_2\\
\dfrac{\tau^4}{8}\mathbf I_2&
\dfrac{\tau^3}{3}\mathbf I_2&
\dfrac{\tau^2}{2}\mathbf I_2\\
\dfrac{\tau^3}{6}\mathbf I_2&
\dfrac{\tau^2}{2}\mathbf I_2&
\tau\mathbf I_2
\end{bmatrix},
\label{eq:q2-process-covariance}
\end{equation}

其中 \(q\) 为 jerk 谱密度。令
\(\mathbf H=[\mathbf I_2\ \mathbf 0\ \mathbf 0]\)。
方式 2 的事件时刻按式\eqref{eq:q2-time-map}校正，位置观测减去
\(\widehat{\mathbf b}\)；两路原始事件按校正时刻合并排序，不将插值点作为额外观测。

对第 \(k\) 个事件，Kalman 预测与创新为

\begin{equation}
\begin{aligned}
\widehat{\mathbf x}_{k|k-1}
&=\mathbf F_k\widehat{\mathbf x}_{k-1|k-1},\\
\mathbf P_{k|k-1}
&=\mathbf F_k\mathbf P_{k-1|k-1}\mathbf F_k^\mathsf T+\mathbf Q_k,\\
\boldsymbol\nu_k
&=\mathbf y_k-\mathbf H\widehat{\mathbf x}_{k|k-1},\\
\mathbf S_{0,k}
&=\mathbf H\mathbf P_{k|k-1}\mathbf H^\mathsf T+\mathbf R_{m(k)}.
\end{aligned}
\label{eq:q2-kalman-prediction}
\end{equation}

定义门控前归一化创新平方

\begin{equation}
\gamma_k=\boldsymbol\nu_k^\mathsf T
\mathbf S_{0,k}^{-1}\boldsymbol\nu_k.
\label{eq:q2-nis}
\end{equation}

设门限 \(g_\alpha=\chi^2_{2,\alpha}\)，按

\begin{equation}
\lambda_k=\max\!\left(1,\frac{\gamma_k}{g_\alpha}\right),
\qquad
\widetilde{\mathbf R}_k=\lambda_k\mathbf R_{m(k)}
\label{eq:q2-robust-gate}
\end{equation}

放大异常观测的协方差。随后计算

\begin{equation}
\begin{aligned}
\mathbf S_k
&=\mathbf H\mathbf P_{k|k-1}\mathbf H^\mathsf T
+\widetilde{\mathbf R}_k,\\
\mathbf K_k
&=\mathbf P_{k|k-1}\mathbf H^\mathsf T\mathbf S_k^{-1},\\
\widehat{\mathbf x}_{k|k}
&=\widehat{\mathbf x}_{k|k-1}+\mathbf K_k\boldsymbol\nu_k,\\
\mathbf P_{k|k}
&=(\mathbf I-\mathbf K_k\mathbf H)\mathbf P_{k|k-1}.
\end{aligned}
\label{eq:q2-kalman-update}
\end{equation}

\noindent\textbf{（6）过程噪声选择与 RTS 平滑。}

过程噪声调参阶段关闭式\eqref{eq:q2-robust-gate}中的协方差膨胀。令

\begin{equation}
\mathbf e_k=
\operatorname{chol}(\mathbf S_{0,k})^{-1}\boldsymbol\nu_k,
\label{eq:q2-whitened-innovation}
\end{equation}

并在候选集 \(\mathcal Q\) 上最小化

\begin{equation}
S(q)=
\left|\overline{\gamma}(q)-2\right|
+2\,\overline{\left|\rho_1(\mathbf e;q)\right|},
\qquad
\widehat q=\arg\min_{q\in\mathcal Q}S(q),
\label{eq:q2-process-score}
\end{equation}

使二维 NIS 均值接近自由度 2，同时降低白化创新的一阶相关。
固定 \(\widehat q\) 后运行稳健门控和完整前向滤波。

离线融合采用 RTS 后向平滑。对 \(k=K-1,\ldots,1\)，有

\begin{equation}
\begin{aligned}
\mathbf G_k
&=\mathbf P_{k|k}\mathbf F_{k+1}^\mathsf T
\mathbf P_{k+1|k}^{-1},\\
\widehat{\mathbf x}_{k|K}
&=\widehat{\mathbf x}_{k|k}
+\mathbf G_k
\left(\widehat{\mathbf x}_{k+1|K}
-\widehat{\mathbf x}_{k+1|k}\right),\\
\mathbf P_{k|K}
&=\mathbf P_{k|k}
+\mathbf G_k
\left(\mathbf P_{k+1|K}-\mathbf P_{k+1|k}\right)
\mathbf G_k^\mathsf T.
\end{aligned}
\label{eq:q2-rts}
\end{equation}

最后仅在两路数据的公共支持域内，以
\(\Delta t_{\mathrm{out}}=0.1\,\mathrm s\) 传播平滑状态，生成 10 Hz
的位置、速度与加速度序列。

\subsubsection{模型二的求解}\label{sec:q2-solution}

\begin{pseudocodeblock}{算法2\quad 稳健时空标定与异步融合}
{\heiti 输入：}两路含噪定位序列
\(\mathcal Z_1,\mathcal Z_2\)，粗搜索集 \(\mathcal H\)，
精搜索半宽 \(L\)，Huber 阈值 \(c\)，IRLS 次数 \(M\)，
过程噪声候选集 \(\mathcal Q\)，门控概率 \(\alpha\)。\par
{\heiti 输出：}时间偏差 \(\widehat{\Delta t}\)、相对空间偏差
\(\widehat{\mathbf b}\) 及 10 Hz 融合状态序列。\par
\begin{enumerate}[label={\heiti 步骤~\arabic*:},leftmargin=5.5em,itemsep=1pt,topsep=3pt]
  \item 构造两路连续位置函数及速度特征，按式\eqref{eq:q2-ncc}计算
  \(C(h)\)，由式\eqref{eq:q2-coarse-offset}得到 \(h_c\)；
  \item 令 \(\mathcal D=[h_c-L,h_c+L]\)，按式\eqref{eq:q2-fixed-grid}
  构造固定公共评价网格，并由式\eqref{eq:q2-robust-scale}确定固定尺度；
  \item 对每个 \(\delta\in\mathcal D\)，迭代 \(M\) 次
  式\eqref{eq:q2-irls-bias}，计算剖面目标 \(J_H(\delta)\)；
  \item 连续最小化式\eqref{eq:q2-huber-objective}，按
  式\eqref{eq:q2-calibration-estimator}得到
  \(\widehat{\Delta t}\) 与 \(\widehat{\mathbf b}\)；
  \item 由式\eqref{eq:q2-third-difference}和
  式\eqref{eq:q2-measurement-covariance}估计
  \(\widehat{\mathbf R}_1,\widehat{\mathbf R}_2\)；
  \item 关闭门控反馈，对 \(q\in\mathcal Q\) 运行前向滤波，按
  式\eqref{eq:q2-process-score}选择 \(\widehat q\)；
  \item 按校正时刻合并原始事件，依次执行
  式\eqref{eq:q2-kalman-prediction}--\eqref{eq:q2-kalman-update}；
  \item 执行式\eqref{eq:q2-rts}的后向平滑，在公共支持域内按
  0.1 s 间隔输出融合状态序列。
\end{enumerate}
\end{pseudocodeblock}

噪声刻画、时空联合标定和异步状态估计之间的关系如
图\ref{fig:q2-framework}所示。
\begin{figure}[htbp]
  \centering
  \includegraphics[width=.98\textwidth]{q2/method_flow.pdf}
  \caption{问题二稳健标定与异步融合流程}
  \label{fig:q2-framework}
\end{figure}

\subsubsection{结果分析}\label{sec:q2-results}

\noindent\textbf{（1）时间与空间偏差估计。}

两种定位方式分别包含 3000 个和 3948 个有效观测。归一化互相关得到粗时间
偏差 \(h_c=52.72\,\mathrm s\)，据此取精搜索区间
\([42.72,62.72]\,\mathrm s\)。固定公共评价网格含 6897 个 10 Hz
时刻，Huber 阈值取 \(c=1.345\)，逐轴 IRLS 迭代 4 次。

图\ref{fig:q2-objective}给出了精搜索区间内的目标函数剖面。Huber
剖面目标在区间内部取得唯一最低点，连续精修值为

\begin{equation}
\widehat{\Delta t}=50.171732\ \mathrm s,
\qquad
t_{2,\mathrm{ref}}=t_2-50.171732\ \mathrm s.
\label{eq:q2-calibration-result}
\end{equation}

\begin{figure}[htbp]
  \centering
  \includegraphics[width=.84\textwidth]{q2/objective_profile.pdf}
  \caption{问题二时空标定目标函数剖面}
  \label{fig:q2-objective}
\end{figure}

将式\eqref{eq:q2-calibration-result}代入逐轴 Huber 估计，得到

\begin{equation}
\widehat{\mathbf b}
=(3.466998,-1.833200)^\mathsf T\ \mathrm m,
\qquad
\|\widehat{\mathbf b}\|_2=3.921824\ \mathrm m.
\label{eq:q2-bias-result}
\end{equation}

该估计表明，方式 2 相对于方式 1 在 \(x\) 方向存在正向平移，在 \(y\)
方向存在负向平移。由于缺少外部真值，\(\widehat{\mathbf b}\)
仅表示两设备的相对系统偏差，不用于判定单个设备的绝对误差。

\noindent\textbf{（2）时空校正效果。}

表\ref{tab:q2-correction}列出校正前后两设备位置差的统计量。

固定网格上的 Huber 标定 RMSE 为 \(1.45262\,\mathrm m\)。
在校正后的 10 Hz 重叠网格上重新计算，两设备位置差的二维 RMSE 为
\(1.45918\,\mathrm m\)。两者的评价网格不同，均表示含噪观测之间的差异，
不是融合轨迹相对于真实轨迹的绝对误差。

\begin{table}[htbp]
\centering
\caption{问题二时空校正前后的两设备位置差}
\label{tab:q2-correction}
\begin{tabular}{lcc}
\toprule
统计量 & 校正前 & 校正后\\
\midrule
\(\overline{\Delta x}\) / m & \(3.47216\) & \(0.00516\)\\
\(\overline{\Delta y}\) / m & \(-1.83035\) & \(0.00286\)\\
\bottomrule
\end{tabular}
\end{table}

校正前后的二维位置关系见图\ref{fig:q2-calibration}。校正前两条轨迹存在明显
平移，按式\eqref{eq:q2-calibration-result}和式\eqref{eq:q2-bias-result}
同时校正时间与空间偏差后，两组含噪观测围绕同一运动轨迹分布。

\begin{figure}[htbp]
  \centering
  \includegraphics[width=.90\textwidth]{q2/calibration_comparison.pdf}
  \caption{问题二时空校正前后的两设备二维位置}
  \label{fig:q2-calibration}
\end{figure}

\noindent\textbf{（3）噪声参数与异步融合。}

三阶差分的 MAD 稳健估计给出

\begin{equation}
\widehat{\mathbf R}_1
=\operatorname{diag}(0.68991,0.71228)\ \mathrm{m}^2,
\qquad
\widehat{\mathbf R}_2
=\operatorname{diag}(0.65003,0.77553)\ \mathrm{m}^2.
\label{eq:q2-measurement-noise-result}
\end{equation}

过程噪声候选集取
\(\mathcal Q=[10^{-11},10^2]\) 上的对数网格。式\eqref{eq:q2-process-score}
在

\begin{equation}
\widehat q=3.162\times10^{-8}
\label{eq:q2-process-noise-result}
\end{equation}

处取得最小值，该值位于候选网格内部。固定 \(\widehat q\) 后采用
\(\alpha=0.99\) 的 NIS 门控。门控前 NIS 均值为 1.9647、95\%
分位为 5.7980，门控后 NIS 均值为 1.9481；白化创新一阶相关绝对值均值为
0.0153，共有 0.892\% 的观测触发协方差放大。

\noindent\textbf{（4）10 Hz 融合轨迹。}

异步 Kalman 滤波与 RTS 平滑后的公共输出区间为
\(102.0\sim951.5\,\mathrm s\)。按 0.1 s 间隔共生成 8496 个状态点，
每个状态同时包含位置、速度和加速度，全部分量均为有限值。
图\ref{fig:q2-trajectory}给出了 10 Hz 融合位置轨迹；起点与终点标记分别对应
公共支持区间的首末时刻。

\begin{figure}[H]
  \centering
  \includegraphics[width=.80\textwidth]{q2/fused_trajectory_10hz.pdf}
  \caption{问题二异步 Kalman/RTS 的 10 Hz 融合轨迹}
  \label{fig:q2-trajectory}
\end{figure}

\subsection{模型三的建立与求解：偏差检验驱动的选择性融合}\label{sec:model-q3}

\subsubsection{模型三的建立}\label{sec:q3-establishment}

\noindent\textbf{（1）实际观测的时间配准。}

设机器人在参考时间轴上的真实轨迹为 \(\mathbf p^\star(t)\)。以方式 1
的时钟为参考，问题三的两组实际定位观测满足

\begin{equation}
\begin{aligned}
\mathbf z_i^{(1)}
&=\mathbf p^\star\!\left(t_i^{(1)}\right)
+\boldsymbol\varepsilon_i^{(1)},\\
\mathbf z_j^{(2)}
&=\mathbf p^\star\!\left(t_j^{(2)}-\Delta t\right)
+\mathbf b(t_j^{(2)})
+\boldsymbol\varepsilon_j^{(2)},
\end{aligned}
\label{eq:q3-observation}
\end{equation}

其中，\(\Delta t\) 为固定时间偏差，\(\mathbf b(t)\) 为待判定的相对系统偏差。
方式 2 的设备时刻映射到参考时间轴为

\begin{equation}
t_{2,\mathrm{ref}}=t_2-\Delta t.
\label{eq:q3-time-map}
\end{equation}

由于两组设备时间区间不直接重叠，先定义各自的经过时间

\begin{equation}
s_i^{(m)}=t_i^{(m)}-t_1^{(m)},\qquad m=1,2,
\label{eq:q3-elapsed-time}
\end{equation}

再对中心化二维位置特征计算归一化互相关，得到粗时间偏差 \(h_c\)。
以 \(h_c\) 为中心构造精搜索域 \(\mathcal D\)，并沿用
式\eqref{eq:q2-fixed-grid}--\eqref{eq:q2-calibration-estimator}
的固定公共网格和 Huber 剖面目标，得到

\begin{equation}
\widehat{\Delta t}
=\arg\min_{\delta\in\mathcal D}J_H(\delta).
\label{eq:q3-time-estimator}
\end{equation}

时间对齐后，在公共评价网格 \(\{t_k\}_{k=1}^{N}\) 上构造未作空间校正的
相对残差

\begin{equation}
\mathbf d_k=
\widetilde{\mathbf z}_2(t_k+\widehat{\Delta t})
-\widetilde{\mathbf z}_1(t_k).
\label{eq:q3-aligned-residual}
\end{equation}

没有外部真值时，\(\mathbf d_k\) 只能用于辨识两设备之间的相对空间差，
不能判断误差应归因于哪一设备。

\noindent\textbf{（2）HAC-Wald 偏差检验。}

检验假设为

\begin{equation}
H_0:\mathbb E(\mathbf d_k)=\mathbf 0,
\qquad
H_1:\mathbb E(\mathbf d_k)\ne\mathbf 0.
\label{eq:q3-hypothesis}
\end{equation}

记样本均值及滞后 \(r\) 阶样本自协方差为

\begin{equation}
\overline{\mathbf d}
=\frac{1}{N}\sum_{k=1}^{N}\mathbf d_k,
\qquad
\widehat{\boldsymbol\Gamma}_r
=\frac{1}{N}\sum_{k=r+1}^{N}
(\mathbf d_k-\overline{\mathbf d})
(\mathbf d_{k-r}-\overline{\mathbf d})^\mathsf T.
\label{eq:q3-autocovariance}
\end{equation}

采用 Newey--West 核构造样本均值的 HAC 协方差估计
\cite{neweywest1987}：

\begin{equation}
\widehat{\boldsymbol\Sigma}_{\overline d}
=\frac{1}{N}
\left[
\widehat{\boldsymbol\Gamma}_0+
\sum_{r=1}^{L}
\left(1-\frac{r}{L+1}\right)
\left(
\widehat{\boldsymbol\Gamma}_r+
\widehat{\boldsymbol\Gamma}_r^\mathsf T
\right)
\right],
\label{eq:q3-hac-covariance}
\end{equation}

其中 \(L\) 为 HAC 截断阶数。Wald 统计量为

\begin{equation}
W=\overline{\mathbf d}^{\,\mathsf T}
\widehat{\boldsymbol\Sigma}_{\overline d}^{-1}
\overline{\mathbf d},
\qquad
W\mathop{\sim}_{H_0}\chi_2^2.
\label{eq:q3-wald}
\end{equation}

\noindent\textbf{（3）工程效应与移动块 Bootstrap。}

统计显著不必然意味着工程上需要校正。令

\begin{equation}
r_k=\|\mathbf d_k\|_2,\qquad
s_r=1.4826\,
\operatorname{median}_k
\left|r_k-\operatorname{median}_{\ell}r_\ell\right|,
\label{eq:q3-residual-scale}
\end{equation}

定义工程效应指数

\begin{equation}
E_b=\frac{\|\overline{\mathbf d}\|_2}{s_r}.
\label{eq:q3-engineering-effect}
\end{equation}

设统计显著性水平为 \(\alpha\)，工程效应阈值为 \(\eta\)，偏差启用判据为

\begin{equation}
I_b=
\mathbf 1\!\left\{
p_{\mathrm{HAC}}<\alpha
\ \land\
E_b\ge\eta
\right\}.
\label{eq:q3-bias-decision}
\end{equation}

为避免依赖渐近分布，同时保留残差的局部时序结构，采用移动块
Bootstrap\cite{kunsch1989}。给定块长 \(\ell\)，构造首尾相接的残差块

\begin{equation}
\mathcal B_j=
\left(
\mathbf d_j,\mathbf d_{j+1},\ldots,\mathbf d_{j+\ell-1}
\right),
\label{eq:q3-bootstrap-block}
\end{equation}

每次有放回抽取若干完整块并截取前 \(N\) 个残差，计算重抽样均值
\(\overline{\mathbf d}^{*(b)}\)。重复 \(B\) 次后，分别取两个坐标方向的
\(\alpha/2\) 与 \(1-\alpha/2\) 分位数，得到逐轴置信区间

\begin{equation}
\mathrm{CI}_r=
\left[
Q_{\alpha/2}
\!\left(\overline d_r^{\,*}\right),
Q_{1-\alpha/2}
\!\left(\overline d_r^{\,*}\right)
\right],
\qquad r\in\{x,y\}.
\label{eq:q3-bootstrap-ci-estimator}
\end{equation}

\noindent\textbf{（4）偏差趋势与状态维度选择。}

为了区分固定偏差与缓慢漂移，对每个坐标拟合

\begin{equation}
d_{kr}=\beta_{0r}+\beta_{1r}t_k+u_{kr},
\qquad r\in\{x,y\},
\label{eq:q3-trend-model}
\end{equation}

并使用 HAC 协方差检验 \(H_0:\beta_{1r}=0\)。令 \(I_{\mathrm{tr}}\)
表示至少一个坐标方向检出显著趋势，则状态模型按

\begin{equation}
\mathcal M=
\begin{cases}
\mathcal M_0,& I_b=0,\\
\mathcal M_{\mathrm c},& I_b=1,\ I_{\mathrm{tr}}=0,\\
\mathcal M_{\mathrm{rw}},& I_b=1,\ I_{\mathrm{tr}}=1
\end{cases}
\label{eq:q3-model-selection}
\end{equation}

选择。其中，\(\mathcal M_0\) 为无偏差状态，
\(\mathcal M_{\mathrm c}\) 为常量偏差状态，
\(\mathcal M_{\mathrm{rw}}\) 为随机游走偏差状态。

无偏差状态与问题二一致：

\begin{equation}
\mathbf x_k^{(0)}
=\left(
\mathbf p_k^\mathsf T,
\mathbf v_k^\mathsf T,
\mathbf a_k^\mathsf T
\right)^\mathsf T.
\label{eq:q3-state-no-bias}
\end{equation}

启用偏差时扩展为

\begin{equation}
\mathbf x_k^{(b)}
=\left(
\mathbf p_k^\mathsf T,
\mathbf v_k^\mathsf T,
\mathbf a_k^\mathsf T,
\mathbf b_k^\mathsf T
\right)^\mathsf T.
\label{eq:q3-state-with-bias}
\end{equation}

扩展状态的测量矩阵为

\begin{equation}
\mathbf H_1=
\begin{bmatrix}
\mathbf I_2&\mathbf 0&\mathbf 0&\mathbf 0
\end{bmatrix},
\qquad
\mathbf H_2=
\begin{bmatrix}
\mathbf I_2&\mathbf 0&\mathbf 0&\mathbf I_2
\end{bmatrix}.
\label{eq:q3-bias-observation-matrix}
\end{equation}

常量偏差模型满足
\(\mathbf b_k=\mathbf b_{k-1}\)；随机游走模型满足

\begin{equation}
\mathbf b_k=\mathbf b_{k-1}+\boldsymbol\xi_k,
\qquad
\boldsymbol\xi_k\sim
\mathcal N(\mathbf 0,q_b\tau_k\mathbf I_2).
\label{eq:q3-bias-random-walk}
\end{equation}

\noindent\textbf{（5）选择性异步融合。}

两设备测量协方差仍由式\eqref{eq:q2-third-difference}和
式\eqref{eq:q2-measurement-covariance}估计。将方式 2 的事件时刻按
式\eqref{eq:q3-time-map}校正后，与方式 1 的原始事件合并排序。
根据式\eqref{eq:q3-model-selection}选择状态维度，再执行
式\eqref{eq:q2-kalman-prediction}--\eqref{eq:q2-rts}的异步 Kalman
滤波和 RTS 平滑。

过程噪声调参阶段关闭门控反馈。对白化创新
\(\mathbf e_k=\operatorname{chol}(\mathbf S_{0,k})^{-1}\boldsymbol\nu_k\)，
在候选集 \(\mathcal Q\) 上最小化

\begin{equation}
S(q)=
\left|\overline{\gamma}(q)-2\right|
+2\,\overline{\left|\rho_1(\mathbf e;q)\right|},
\qquad
\widehat q=\arg\min_{q\in\mathcal Q}S(q).
\label{eq:q3-process-score}
\end{equation}

固定 \(\widehat q\) 后运行 NIS 稳健门控。对于 10 Hz 非事件时刻，从前一
RTS 平滑事件状态按状态方程传播。由输出协方差的位置子矩阵
\(\mathbf P_{p,k}\) 定义二维合成位置标准差

\begin{equation}
\sigma_{p,k}=
\sqrt{\operatorname{tr}(\mathbf P_{p,k})}.
\label{eq:q3-position-std}
\end{equation}

\subsubsection{模型三的求解}\label{sec:q3-solution}

\begin{pseudocodeblock}{算法3\quad 偏差检验驱动的选择性融合}
{\heiti 输入：}两路实际定位序列
\(\mathcal Z_1,\mathcal Z_2\)，粗搜索集 \(\mathcal H\)，
精搜索半宽 \(L_t\)，显著性水平 \(\alpha\)，工程阈值 \(\eta\)，
HAC 阶数 \(L\)，Bootstrap 次数 \(B\) 与块长 \(\ell\)，
过程噪声候选集 \(\mathcal Q\)。\par
{\heiti 输出：}时间偏差 \(\widehat{\Delta t}\)、偏差判定
\(I_b\)、所选状态模型 \(\mathcal M\) 及 10 Hz 融合状态。\par
\begin{enumerate}[label={\heiti 步骤~\arabic*:},leftmargin=5.5em,itemsep=1pt,topsep=3pt]
  \item 按式\eqref{eq:q3-elapsed-time}重置两路时间，利用中心化位置特征的
  归一化互相关获得粗时间偏差 \(h_c\)；
  \item 在 \([h_c-L_t,h_c+L_t]\) 内最小化固定网格 Huber 剖面目标，
  由式\eqref{eq:q3-time-estimator}得到 \(\widehat{\Delta t}\)；
  \item 按式\eqref{eq:q3-aligned-residual}构造同步残差，计算
  \(\overline{\mathbf d}\)、HAC 协方差、Wald 统计量和工程效应指数；
  \item 重复 \(B\) 次移动块 Bootstrap，按
  式\eqref{eq:q3-bootstrap-ci-estimator}构造逐轴置信区间；
  \item 拟合式\eqref{eq:q3-trend-model}并执行 HAC 趋势检验，按
  式\eqref{eq:q3-bias-decision}和式\eqref{eq:q3-model-selection}
  选择状态模型；
  \item 由位置三阶差分估计两设备测量协方差，关闭门控反馈并按
  式\eqref{eq:q3-process-score}选择 \(\widehat q\)；
  \item 按校正事件时刻执行所选模型的异步 Kalman 滤波与 NIS 稳健门控，
  再执行 RTS 后向平滑；
  \item 在公共支持域内按 0.1 s 间隔传播状态，输出位置、速度、加速度及
  式\eqref{eq:q3-position-std}定义的模型内位置标准差。
\end{enumerate}
\end{pseudocodeblock}

时间配准、偏差检验、状态模型选择和事件融合之间的关系如
图\ref{fig:q3-framework}所示。
\begin{figure}[htbp]
  \centering
  \includegraphics[width=.98\textwidth]{q3/method_flow.pdf}
  \caption{问题三偏差检验与选择性融合流程}
  \label{fig:q3-framework}
\end{figure}

\subsubsection{结果分析}\label{sec:q3-results}

\noindent\textbf{（1）时间配准结果。}

两种定位方式分别包含 1381 个和 1621 个有效观测。经过时间上的中心化位置
互相关得到粗时间偏差

\begin{equation}
h_c=-368.82\ \mathrm s.
\label{eq:q3-coarse-result}
\end{equation}

以 \(h_c\) 为中心取 10 s 精搜索半宽，搜索区间为
\([-378.82,-358.82]\,\mathrm s\)。固定公共评价网格含 2910 个时刻，
Huber 剖面连续精修得到

\begin{equation}
\widehat{\Delta t}=-367.877619\ \mathrm s,
\qquad
t_{2,\mathrm{ref}}
=t_2-(-367.877619)
=t_2+367.877619\ \mathrm s.
\label{eq:q3-time-offset}
\end{equation}

因此，方式 2 的设备时钟相对于方式 1 参考时钟落后约
\(367.88\,\mathrm s\)，映射时需将方式 2 时间戳整体增加相应数值。

固定网格 Huber 剖面的伴随位置中心为

\begin{equation}
\widehat{\mathbf b}_{\mathrm{prof}}
=(-0.151775,-0.150330)^\mathsf T\ \mathrm m.
\label{eq:q3-profile-bias-result}
\end{equation}

相应标定 RMSE 为 \(6.36335\,\mathrm m\)。该位置中心仅作为时间剖面目标中的
伴随参数，系统偏差是否启用由后续统计检验决定。

\noindent\textbf{（2）系统偏差统计判定。}

时间配准后用于偏差推断的残差序列含 3000 个 10 Hz 时刻，其普通均值为

\begin{equation}
\overline{\mathbf d}
=(-0.133731,-0.157069)^\mathsf T\ \mathrm m,
\qquad
\|\overline{\mathbf d}\|_2=0.206288\ \mathrm m.
\label{eq:q3-mean-residual}
\end{equation}

HAC 截断阶数取 \(L=8\)，得到样本均值协方差

\begin{equation}
\widehat{\boldsymbol\Sigma}_{\overline d}
=
\begin{bmatrix}
0.020304&0.001455\\
0.001455&0.018041
\end{bmatrix}
\mathrm m^2.
\label{eq:q3-hac-result}
\end{equation}

取 \(\alpha=0.05\)、工程阈值 \(\eta=0.25\)，检验结果见
表\ref{tab:q3-decision}。Wald 检验未拒绝零相对偏差，工程效应指数也未达到
启用偏差状态的阈值。

\begin{table}[htbp]
\centering
\caption{问题三相对系统偏差的统计与工程判定}
\label{tab:q3-decision}
\begin{tabular}{lcc}
\toprule
判定量 & 计算结果 & 判定阈值\\
\midrule
HAC-Wald 统计量 \(W\) & \(2.0935\) & \(\chi^2_{2,0.95}\)\\
\(p_{\mathrm{HAC}}\) & \(0.3511\) & \(0.05\)\\
稳健噪声尺度 \(s_r\) / m & \(3.0300\) & --\\
工程效应指数 \(E_b\) & \(0.0681\) & \(0.25\)\\
\bottomrule
\end{tabular}
\end{table}

移动块 Bootstrap 取 \(B=2000\)、块长 \(\ell=14\)，所得 95\%
置信区间为

\begin{equation}
b_x\in[-0.413289,0.138470]\ \mathrm m,
\qquad
b_y\in[-0.428052,0.119683]\ \mathrm m.
\label{eq:q3-bootstrap-ci}
\end{equation}

两轴区间均覆盖零。图\ref{fig:q3-bias-ci}同时给出了样本均值和移动块
Bootstrap 区间，零水平线均穿过两个坐标方向的置信区间。

\begin{figure}[htbp]
  \centering
  \includegraphics[width=.76\textwidth]{q3/bias_confidence_interval.pdf}
  \caption{问题三相对偏差估计及 95\% 移动块 Bootstrap 置信区间}
  \label{fig:q3-bias-ci}
\end{figure}

x、y 方向的 HAC 趋势斜率分别为
\(-1.4675\times10^{-3}\,\mathrm{m/s}\) 和
\(1.4894\times10^{-4}\,\mathrm{m/s}\)，对应 \(p\) 值为
0.3783 和 0.9208，均未检出显著漂移。图\ref{fig:q3-bias-series}
显示残差存在明显的局部时序相关，因此偏差均值的协方差不能按独立样本估计；
两条滚动均值曲线则未呈现贯穿观测区间的单调漂移。

\begin{figure}[htbp]
  \centering
  \includegraphics[width=.86\textwidth]{q3/bias_time_series.pdf}
  \caption{问题三时间配准后的两坐标相对残差时序}
  \label{fig:q3-bias-series}
\end{figure}

综合 \(p_{\mathrm{HAC}}\)、\(E_b\)、Bootstrap 区间和趋势检验，在给定
显著性水平与工程阈值下，未发现具有统计与工程意义的相对系统偏差，
即 \(I_b=0\)，选择六维无偏差状态模型 \(\mathcal M_0\)。
该结论表示现有样本不足以支持启用偏差状态，不等同于证明真实偏差严格为零。

\noindent\textbf{（3）噪声参数与异步融合。}

位置三阶差分的 MAD 估计给出

\begin{equation}
\widehat{\mathbf R}_1
=\operatorname{diag}(16.0645,18.2067)\ \mathrm{m}^2,
\qquad
\widehat{\mathbf R}_2
=\operatorname{diag}(7.2399,8.9392)\ \mathrm{m}^2.
\label{eq:q3-measurement-noise-result}
\end{equation}

过程噪声候选集取 \([10^{-6},10^3]\) 上的对数网格。
式\eqref{eq:q3-process-score}在

\begin{equation}
\widehat q=3.1623
\label{eq:q3-process-noise-result}
\end{equation}

处取得最小值，该值位于候选区间内部。图\ref{fig:q3-tuning}给出了
过程噪声谱密度与 NIS/白化创新综合得分的关系。

\begin{figure}[htbp]
  \centering
  \includegraphics[width=.80\textwidth]{q3/process_noise_tuning.pdf}
  \caption{问题三基于门控前创新统计的过程噪声选择}
  \label{fig:q3-tuning}
\end{figure}

固定 \(\widehat q\) 后采用 99\% NIS 门控。门控前平均 NIS 为
2.0127、95\% 分位为 5.9590，门控后平均 NIS 为 1.9957；
白化创新平均绝对一阶相关为 0.0274，共有 1.033\% 的观测触发协方差放大。

\noindent\textbf{（4）10 Hz 融合状态与不确定性。}

融合轨迹的公共输出区间为
\(445.0076\sim814.0076\,\mathrm s\)，按 0.1 s 间隔共生成
3691 个状态点。位置、速度、加速度和协方差派生量均为有限值；
10 Hz 非事件时刻从前一平滑事件节点传播的最大时间间隔为 0.208 s。

二维合成位置标准差的中位数为 \(0.9606\,\mathrm m\)，最大值为
\(3.4593\,\mathrm m\)。图\ref{fig:q3-trajectory}以颜色表示
\(\sigma_{p,k}\)，并标记公共输出区间的起点和终点。该标准差只反映所建
状态空间模型内的不确定性，不等同于融合轨迹相对于未知真实轨迹的绝对误差。

\begin{figure}[H]
  \centering
  \includegraphics[width=.72\textwidth]{q3/fused_trajectory_uncertainty.pdf}
  \caption{问题三 10 Hz 融合轨迹及模型内位置标准差}
  \label{fig:q3-trajectory}
\end{figure}

\subsection{模型四的建立与求解：固定轨迹任务优化}\label{sec:model-q4}

\subsubsection{模型四的建立}\label{sec:q4-establishment}

\noindent\textbf{（1）轨迹状态、目标距离与拍照方向角。}

问题四不改变问题三所得轨迹，只在固定轨迹上选择任务和时刻。设机器人状态为
\(\mathbf p(t)=(x(t),y(t))^{\mathsf T}\)、速度模 \(v(t)\)、加速度模
\(a(t)\)，目标 \(j\) 的坐标为 \(\mathbf g_j\)，则距离为

\begin{equation}
d_j(t)=\|\mathbf p(t)-\mathbf g_j\|_2.
\label{eq:q4-distance}
\end{equation}

射击任务在执行前连续 1.5 s 内须满足

\begin{equation}
5\le d_j(t)\le30,\qquad v(t)\le2,\qquad a(t)\le1.5;
\label{eq:q4-shooting-constraints}
\end{equation}

每个射击目标至多执行一次，单次命中概率为 0.85。拍照任务在执行前连续 0.5
s 内须满足

\begin{equation}
10\le d_j(t)\le40,\qquad v(t)\le1.5,\qquad a(t)\le1.5.
\label{eq:q4-photo-constraints}
\end{equation}

机器人相对拍照目标的方向角定义为

\begin{equation}
\theta_j(t)=\operatorname{atan2}(y(t)-y_j,x(t)-x_j)\pmod {360^\circ},
\label{eq:q4-photo-angle}
\end{equation}

同一目标两次拍照的圆周角距离不得小于
\(60^\circ\)。题面未明示单机器人能否并行准备，也未规定同一射击目标是否可重复计数。为使``完成目标''具有明确工程含义，本文作两条建模假设：机器人只有一套任务执行资源，任意两项已选任务的准备/执行闭区间不得相交，并额外留出
0.01 s
数值间隔；每个射击目标至多计一次完成。结果总数还受模板九行容量限制。这两条属于模型假设而非题面原文，后续约束与最优性结论均以此为边界。

\noindent\textbf{（2）连续可行窗口与候选任务。}

先在正式 10 Hz
状态上，对每个可能执行时刻向前检查完整准备窗的全部采样点。对距离下界、距离上界、速度和加速度分别归一化，候选
\(c\) 的最小裕度定义为

\begin{equation}
m_c=\min_{t\in[t_c-L_c,t_c]}
\left\{
\frac{d-d_{\min}}{d_{\max}-d_{\min}},
\frac{d_{\max}-d}{d_{\max}-d_{\min}},
\frac{v_{\max}-v}{v_{\max}},
\frac{a_{\max}-a}{a_{\max}}
\right\}.
\label{eq:q4-margin}
\end{equation}

在每个连续可行段内保留定时间隔点、局部高裕度点；拍照任务另按 30°
方向角箱保留高裕度点。随后用 PCHIP 连续状态在 0.01 s
网格重新检查完整准备区间，剔除任何负裕度候选。共得到 1255
个候选，其中射击 595 个、拍照 660 个。

\noindent\textbf{（3）资源互斥、射击唯一性与词典序MILP。}

令 \(x_c\in\{0,1\}\) 表示是否选择候选 \(c\)。对时间冲突候选对
\((c,d)\in\mathcal T\)、同一射击目标候选集
\(\mathcal S_j\)，以及同一拍照目标且角差不足 60° 的候选对
\((c,d)\in\mathcal A\)，约束为

\begin{equation}
x_c+x_d\le1\quad((c,d)\in\mathcal T\cup\mathcal A),
\qquad
\sum_{c\in\mathcal S_j}x_c\le1,
\qquad
\sum_cx_c\le9.
\label{eq:q4-logical-constraints}
\end{equation}

第一级目标为

\begin{equation}
\max\sum_cx_c.
\label{eq:q4-milp-stage1}
\end{equation}

固定最优任务数 \(N^*\) 后，引入最小裕度 \(z\)，以充分大的 \(M\) 施加

\begin{equation}
z\le m_c+M(1-x_c),\qquad \max z.
\label{eq:q4-milp-stage2}
\end{equation}

最后固定 \(N^*\) 和
\(z^*\)，最大化总裕度；对可同时选择的同目标拍照对增加 AND
辅助变量及按圆周角差归一化的多样性奖励，射击的 0.85
期望命中仅作为很小的同级择优项。三个阶段均由 HiGHS 返回最优状态且 MIP
gap 为
0。求解器的线性优化实现背景参见文献\cite{huangfu2018}，但本文的最优性结论只依据实际求解状态、零 gap 与独立复核。由于九行是不可突破的题面输出上界，找到九项可行任务也证明了任务数量已达全局上界；裕度和多样性的最优性针对经连续验证的压缩候选集。

\subsubsection{模型四的求解}\label{sec:q4-solution}

求解先识别连续可行窗口，再以0.01 s网格细化并压缩候选；随后依次最大化任务数量、最小安全裕度和总裕度/拍照多样性。三个阶段均固定前一阶段最优值，故次级目标不会牺牲任务数量。最终对候选窗口作0.001 s连续复核，并在时间保留两位小数后再次核验。

候选可行性分布见图\ref{fig:q4-candidates}。候选集中于轨迹经过目标邻域且速度、加速度同时满足完整准备窗的区段，这一分布来自物理约束，而不是优化器任意删选。

\begin{figure}[htbp]
  \centering
  \includegraphics[width=.90\textwidth]{q4/candidate_feasibility_map.pdf}
  \caption{模型四经连续准备窗复核后的1255个候选任务分布}
  \label{fig:q4-candidates}
\end{figure}

\subsubsection{算法实现}\label{sec:q4-algorithm}

\begin{pseudocodeblock}{算法4\quad 固定轨迹上的任务调度}
{\heiti 输入：}问题三10 Hz状态、模型内位置标准差、目标坐标及任务约束。\par
{\heiti 输出：}九项任务的目标编号、任务类型、准备时刻与执行时刻。\par
\begin{enumerate}[label={\heiti 步骤~\arabic*:},leftmargin=5.5em,itemsep=1pt,topsep=3pt]
  \item 插值轨迹状态，逐目标检查射击或拍照的完整准备区间，识别连续可行窗口；
  \item 在每个窗口内保留定时间隔点、局部高裕度点及拍照方向角代表点，并在0.01 s网格上复核；
  \item 建立时间冲突、角度冲突、射击唯一性和九行容量约束，依次求解任务数、最小裕度和综合质量三级MILP；
  \item 若任一级未返回最优状态或MIP gap不为0，则终止并报告失败；否则进入连续复核；
  \item 在0.001 s网格上检查所选任务，时刻保留两位小数后再次检查全部准备窗与冲突约束；
  \item 将通过复核的九行结果写入工作簿副本并返回日程，原始模板保持不变。
\end{enumerate}
\end{pseudocodeblock}

\subsubsection{结果分析}\label{sec:q4-results}

正式方案见表\ref{tab:q4-schedule}。输出时刻按模板保留两位小数，并在舍入后再次进行 0.01
s 全准备窗复核。

\begin{table}[htbp]
\centering
\caption{问题四正式九项任务日程}
\label{tab:q4-schedule}
\begin{tabular}{rccrr}
\toprule
序号 & 目标编号 & 任务 & 开始准备时刻/s & 任务执行时刻/s \\
\midrule
1 & S05 & 射击 & 450.31 & 451.81 \\
2 & P04 & 拍照 & 454.61 & 455.11 \\
3 & P09 & 拍照 & 455.21 & 455.71 \\
4 & P12 & 拍照 & 456.01 & 456.51 \\
5 & S04 & 射击 & 464.21 & 465.71 \\
6 & P16 & 拍照 & 466.61 & 467.11 \\
7 & P03 & 拍照 & 493.21 & 493.71 \\
8 & S11 & 射击 & 509.31 & 510.81 \\
9 & S10 & 射击 & 525.51 & 527.01 \\
\bottomrule
\end{tabular}
\end{table}

方案含 4 次射击、5 次拍照，射击期望命中数为

\begin{equation}
E(N_{\mathrm{hit}})=4\times0.85=3.40.
\label{eq:q4-expected-hits}
\end{equation}

五次拍照分属五个不同目标，因此同目标角差约束自动满足。正式方案最小归一化约束裕度为
0.4241；最早可行贪心基线虽同样完成九项，但最小裕度仅 0.1010，MILP
方案约为其 4.20
倍。全部任务在准备区间内的距离、速度和加速度均已通过连续复核。

轨迹、目标位置和所选任务见图\ref{fig:q4-layout}，日程时间线见图\ref{fig:q4-timeline}。任务集中在450.31--527.01 s是固定轨迹与目标几何共同作用的结果；所有准备/执行闭区间互不重叠，最紧相邻任务仍保留0.10 s间隔。

\subsubsection{模型检验与稳健性}\label{sec:q4-validation}

设置三档压力情景，每档使用固定随机种子生成 500 次扰动：中等情景采用 0.5
倍问题三模型位置标准差和 5\% 速度/加速度相对扰动，标准情景采用 1.0
倍位置标准差和 5\% 状态扰动，严重情景采用 1.5 倍位置标准差和 10\%
状态扰动。九项任务逐项和整套方案在三档情景中的可行率均为
100\%。该结果只说明方案在所定义扰动集下具有较大裕度，不被解释为真实任务成功概率；射击命中仍单独按题设
0.85 计算。

最终将九行结果写入 \texttt{result.xlsx} 副本的 B2:E10。A 列序号、H:L
红色说明和范例、其余单元格值及样式均经语义快照核对保持不变；原始模板未被覆盖。
约束裕度见图\ref{fig:q4-margins}。
全部1255个候选的裕度经验分布见图\ref{fig:q4-margin-cdf}。

\begin{figure}[htbp]
  \centering
  \includegraphics[width=.88\textwidth]{q4/trajectory_targets_schedule.pdf}
  \caption{模型四固定轨迹、目标点与九项所选任务}
  \label{fig:q4-layout}
\end{figure}
\begin{figure}[htbp]
  \centering
  \includegraphics[width=.88\textwidth]{q4/optimized_schedule_timeline.pdf}
  \caption{模型四九项任务的准备与执行时间线}
  \label{fig:q4-timeline}
\end{figure}
\begin{figure}[htbp]
  \centering
  \includegraphics[width=.90\textwidth]{q4/constraint_margins.pdf}
  \caption{模型四九项任务的归一化完整准备窗约束裕度}
  \label{fig:q4-margins}
\end{figure}
\begin{figure}[htbp]
  \centering
  \includegraphics[width=.78\textwidth]{paper/candidate_margin_cdf.pdf}
  \caption{模型四射击与拍照候选裕度的经验累积分布}
  \label{fig:q4-margin-cdf}
\end{figure}

